\documentclass[aspectratio=169,11pt]{beamer}
% Shared look for every deck: Metropolis with a progress bar under each frame title.
% Build with Tectonic: `make` in this directory.

\usetheme[progressbar=frametitle,numbering=fraction,sectionpage=progressbar,block=fill]{metropolis}

% Metropolis looks Fira Sans up by system name; load it from the TeX bundle instead,
% Light for body text as Metropolis intends, with lining figures for data.
\usepackage[sfdefault,light,lf]{FiraSans}
\usepackage{booktabs,tabularx}
\usepackage{tikz}
\usetikzlibrary{arrows.meta,positioning,shapes.geometric,calc}

% Metropolis's palette, with the decks' own colour names mapped onto it.
\definecolor{mDarkTeal}{HTML}{23373B}
\definecolor{mDarkBrown}{HTML}{604C38}
\definecolor{mLightBrown}{HTML}{EB811B}
\colorlet{crust}{mDarkBrown}
\definecolor{crumb}{HTML}{FAEEE2}
\colorlet{rye}{mDarkTeal}
\definecolor{flag}{HTML}{B03A2E}
\definecolor{own}{HTML}{2E7D4F}
\definecolor{muted}{HTML}{7A7F80}

% Top-align frames, so the \vfill in \source puts the source line at the bottom.
% A slimmer title bar and narrower margins than Metropolis's defaults leave room
% for the denser slides.
\makeatletter
\beamer@centeredfalse
\setlength{\metropolis@frametitle@padding}{1.4ex}
% A progress bar you can see: under every frame title and on section pages.
\setlength{\metropolis@progressinheadfoot@linewidth}{2pt}
\setlength{\metropolis@progressonsectionpage@linewidth}{1.5pt}
\setlength{\metropolis@titleseparator@linewidth}{1pt}
% Metropolis's title page is \paperheight tall, which overflows a top-aligned frame.
\setbeamertemplate{title page}{%
  \begin{minipage}[b][\textheight]{\textwidth}
    \vfill
    \usebeamertemplate*{title}
    \ifx\insertsubtitle\@empty\else\usebeamertemplate*{subtitle}\fi
    \usebeamertemplate*{title separator}
    \usebeamertemplate*{author}
    \usebeamertemplate*{institute}
    \vfill
  \end{minipage}}
\makeatother
\setbeamersize{text margin left=1.2em,text margin right=1.2em}

\setbeamertemplate{frame footer}{\insertshorttitle}
\setbeamertemplate{footline}{%
  \begin{beamercolorbox}[wd=\paperwidth,ht=2.6ex,dp=1.2ex,leftskip=1.2em,rightskip=1.2em]{footline}%
    \usebeamerfont{page number in head/foot}%
    \usebeamertemplate*{frame footer}\hfill\usebeamertemplate*{frame numbering}%
  \end{beamercolorbox}}
\setbeamerfont{page number in head/foot}{size=\tiny}
\setbeamercolor{footline}{fg=muted}

% The practical point of a slide, for the baker.
\newcommand{\takeaway}[1]{%
  \par\medskip
  \colorbox{crumb}{\parbox{\dimexpr\linewidth-2\fboxsep}{\small\textbf{\color{crust}In the kitchen:} #1}}}

% Small source line, placed at the bottom of a frame.
\newcommand{\source}[1]{%
  \vfill{\color{muted}\tiny Source: #1}}

% Status flags, mirroring the notes' conventions. Filled tags, so they read on the
% dark title bar as well as on the slide body.
\newcommand{\flagtag}[2]{\mbox{\setlength{\fboxsep}{1.2pt}\colorbox{#1}{\color{white}\scriptsize\bfseries\,#2\,}}}
\newcommand{\hypothesis}{\flagtag{flag}{hypothesis}}
\newcommand{\contested}{\flagtag{flag}{contested}}
\newcommand{\ownresult}{\flagtag{own}{own result}}
\newcommand{\untested}{\flagtag{flag}{not yet tested}}

% Source lists: one source per bullet, author and year in bold, optional group
% headings. Use in a [allowframebreaks] frame; long lists continue on "(cont.)" slides.
\newenvironment{sourcelist}{%
  \setbeamerfont{itemize/enumerate body}{size=\scriptsize}%
  \setbeamerfont{itemize item}{size=\scriptsize}%
  \scriptsize\begin{itemize}\setlength{\itemsep}{2pt}}{\end{itemize}}
\newcommand{\src}[2]{\item\textbf{#1}\enspace #2}
% \nobreak keeps a group heading on the same slide as its first source.
\newcommand{\srcgroup}[1]{\par\smallskip{\small\bfseries\color{crust}#1}\par\nobreak\vspace{-2pt}\nobreak}
\setbeamertemplate{frametitle continuation}[from second][(cont.)]

% Flowchart nodes for TikZ diagrams.
\tikzset{
  box/.style={draw=crust,fill=crumb,rounded corners=2pt,align=center,
    font=\scriptsize,inner sep=4pt,text width=#1},
  box/.default=2.6cm,
  arrow/.style={-{Stealth[length=5pt]},draw=crust,thick},
}

\author{Sourdough Science}
\institute{github.com/zrrbite/sourdough-science}
\date{}


% Local additions for this deck only.
% A short inline qualifier, as a filled tag like the status flags.
\newcommand{\caveat}[1]{\flagtag{flag}{#1}}
% A 7 pt size for dense diagrams.
\newcommand{\diagfont}{\fontsize{7}{8.4}\selectfont}
\tikzset{
  stage/.style={draw=crust,fill=crust,text=white,rounded corners=2pt,align=center,
    font=\scriptsize\bfseries,inner sep=3pt,text width=#1,minimum height=0.55cm},
  stage/.default=3.05cm,
  cell/.style={draw=crust!50,fill=crumb,rounded corners=2pt,align=left,
    font=\diagfont,inner sep=3pt,text width=#1,
    execute at begin node={\hyphenpenalty=10000\exhyphenpenalty=10000}},
  cell/.default=3.05cm,
  skipped/.style={draw=muted,dashed,fill=white,text=muted,rounded corners=2pt,
    align=left,font=\diagfont,inner sep=3pt,text width=#1,
    execute at begin node={\hyphenpenalty=10000\exhyphenpenalty=10000}},
  rowlabel/.style={font=\scriptsize\bfseries,text=muted,align=right,text width=1.3cm},
  nohyph/.style={execute at begin node={\hyphenpenalty=10000\exhyphenpenalty=10000}},
  tcell/.style={draw=crust!50,fill=crumb,rounded corners=2pt,align=left,
    font=\scriptsize,inner sep=3pt,text width=3.05cm,
    execute at begin node={\hyphenpenalty=10000\exhyphenpenalty=10000\hbadness=10000}},
}

\title[Starters]{Starters: From Scratch, Dried, and Revived}
\subtitle{The 14-day method, what is happening in the jar,\\and a backup that comes back from a pinch}

\begin{document}

\begin{frame}[plain]\titlepage\end{frame}

% ---------------------------------------------------------------------------
\begin{frame}{What this deck covers}
  \begin{enumerate}
    \item \textbf{The 14-day method}: Full Proof Baking's schedule in grams, and her rules
    \item \textbf{What is happening in the jar}: the early rise is bacteria; the day-5 lull is the hand-over
    \item \textbf{Video, literature and this kitchen}: where the explanations differ; warm, not hot
    \item \textbf{Drying a backup and reviving it}: a pinch, not a spoonful
    \item \textbf{Open tests}: seven experiments the note leaves open
  \end{enumerate}
  \bigskip
  {\small Flags carried over from the notes: \hypothesis\ \untested\ \caveat{not peer-reviewed}\ \caveat{loose analogue}.
  Results from this kitchen are marked \ownresult; everything else is cited at the foot of the slide.}
  \source{notes/starter-from-scratch.md and notes/starter-cheatsheet.md}
\end{frame}

% ===========================================================================
\section{The 14-day method}

\begin{frame}{What you need: a 30\% rye mix, a rubber band and 24°C}
  \begin{itemize}
    \item \textbf{Flour mix, 30\% rye:} 350\,g bread flour + 150\,g dark rye, whisked together. Every feed until day 12 uses it
    \item \textbf{A straight-sided jar} with a rubber band at the level of each fresh feed; a fresh jar on day 2, and whenever it gets messy
    \item \textbf{A scale reading to 0.1\,g}, tared before each ingredient: the feeds go down to 5\,g
    \item \textbf{24°C} (75°F), ``warm room temperature'', for all 14 days
    \item \textbf{Cheesecloth} and a rubber band until day 8, then a loosely fitted lid
  \end{itemize}
  \medskip
  Ratios are \textbf{starter\,:\,water\,:\,flour by weight}, so every feed is at 100\% hydration.
  \takeaway{Rye for its nutrients and enzymes, bread flour for strength. Plain flour also works for the base.}
  \source{Kristen, Full Proof Baking (2019), from her on-screen title cards and narration}
\end{frame}

\begin{frame}{Two weeks in grams: 1:1:1 for a week, then the ratio climbs}
  \footnotesize
  \renewcommand{\arraystretch}{0.95}
  \begin{tabularx}{\linewidth}{@{}c c c X@{}}
    \toprule
    \textbf{Day} & \textbf{Morning, g} & \textbf{Evening, g} & \textbf{What you'll see} \\
    \midrule
    1   & --- & 0 / 30 / 30 & Mix only, rest 24\,h \\
    2   & --- & Stir, clean jar & Dry skin on top: fine \\
    3   & --- & 20 / 20 / 20 & Risen: bacteria, not yeast yet \\
    4   & 20 / 20 / 20 & 16 / 16 / 16 & 3--4×, runny, cheesy smell \\
    \textbf{5} & \textbf{15 / 15 / 15} & \textbf{16 / 16 / 16} & \textbf{Almost no rise, vomit smell: normal} \\
    6--7 & 15 / 15 / 15 & 15 / 15 / 15 & 1.2--2×, milder, first sour notes \\
    8   & 15 / 15 / 15 & 15 / 15 / 15 & \textasciitilde 2×. Save discard, loose lid \\
    9   & 15 / 15 / 15 & 8 / 16 / 16 & 2.5× \\
    10  & 8 / 16 / 16 & 5 / 15 / 15 & 2.5×, falls before the feed \\
    11  & 5 / 20 / 20 & 5 / 25 / 25 & 2.5× \\
    12  & 5 / 25 / 25 & 5 / 20 / 20 & Evening: switch to 10\% rye mix \\
    13  & 5 / 25 / 25 & 5 / 25 / 25 & 2.5× \\
    14  & 8 / 16 / 16 ×2 & 10 / 50 / 50 & Feeds 7\,h apart; three a day now \\
    \bottomrule
  \end{tabularx}
  \source{Full Proof Baking (2019), as tabulated in the cheat sheet. Feeds in grams of starter / water / flour: two a day from day 4, three from day 14}
\end{frame}

\begin{frame}{Feed at peak; it is done when it triples in 4--6\,h at 27°C}
  \small
  \begin{columns}[T]
    \begin{column}{0.5\textwidth}
      \textbf{\color{crust}Feeding}
      \begin{itemize}
        \item Fallen before the feed? Next feed one ratio up: 1:2:2 → 1:3:3 → 1:4:4 → 1:5:5
        \item Fewer than three feeds a day: higher ratios, up to 1:10:10
        \item Day 12: 10\% rye mix = leftover 30\% mix + twice its weight in bread flour, for a ``higher peak, hold peak longer''
        \item \textbf{Discard} if pink, orange or fuzzy
      \end{itemize}
    \end{column}
    \begin{column}{0.46\textwidth}
      \textbf{\color{crust}Done}
      \begin{itemize}
        \item Triples in 4--6\,h on 1:2:2 at \textbf{27°C}, not room temperature
        \item No 40°C+ builds until done. Fridge: 2--4 weeks unfed
      \end{itemize}
      \smallskip
      {\footnotesize
      \begin{tabular}{@{}lll@{}}
        \toprule
        Video, day 14 & 2× at & Peak \\
        \midrule
        1:5:5 & 5\,h & 2.9× at 7.5\,h \\
        1:2:2 & 4\,h & 3× at 7\,h \\
        \bottomrule
      \end{tabular}\par\smallskip
      Not there yet: ``still pretty young''}
    \end{column}
  \end{columns}
  \takeaway{Her closing card: ``Feed at peak! Monitor your starter! Keep it warm!''}
  \source{Full Proof Baking (2019); the day-14 time-lapse, in an oven with the light on at 27°C (80°F)}
\end{frame}

% ===========================================================================
\section{What is happening in the jar}

\begin{frame}{The first two weeks are a succession of microbes}
  \centering
  \begin{tikzpicture}[x=1cm,y=1cm]
    \foreach \i/\head in {0/{Days 1--4\\the first wave},1/{Day 5\\the lull},2/{Days 6--9\\the hand-over},3/{Days 9--14\\building strength}}{
      \node[stage] (h\i) at (2.95+\i*3.35,0) {\head};
    }
    \foreach \i in {0,1,2}{
      \pgfmathtruncatemacro{\j}{\i+1}
      \draw[arrow] (h\i.east) -- (h\j.west);
    }
    \node[rowlabel,anchor=east] at (1.3,-1.7) {In the jar};
    \node[rowlabel,anchor=east] at (1.3,-3.4) {She saw};
    \foreach \x/\t in {2.95/{\textit{Leuconostoc}, \textit{Lactococcus}, \textit{Enterococcus}, \textit{Enterobacter} from the flour},
                       6.30/{Acid stops the first wave; the sourdough organisms are not yet in numbers},
                       9.65/{\textit{Lactobacillus}, \textit{Pediococcus}, \textit{Weissella} take over; yeast the day after pH 3.5},
                       13.00/{Community largely settled; what remains is matching food to appetite}}{
      \node[tcell,minimum height=2.2cm,anchor=north] at (\x,-0.6) {\t};
    }
    \foreach \x/\t in {2.95/{3--4× rise, unpleasant smells},
                       6.30/{``No significant growth!'', the worst smell},
                       9.65/{1.2× → 2×; ``bit of sour'' on day 7},
                       13.00/{Steady 2.5×; the ratio climbs}}{
      \node[tcell,minimum height=0.9cm,anchor=north] at (\x,-2.95) {\t};
    }
  \end{tikzpicture}
  \raggedright
  \takeaway{The early rise that dies and the smell put people off, but they are an ordinary, well-documented succession.}
  \source{The two lab studies, both refreshed daily: Van der Meulen et al.\ (2007) at 30°C; Ercolini et al.\ (2013) 5\,h a day at 25°C, 10°C between. Also Wink (2008a, b); Full Proof Baking (2019)}
\end{frame}

\begin{frame}{Days 1--4: the early rise is bacteria, not yeast}
  \small
  \begin{itemize}
    \item Flour brings mostly \textit{Proteobacteria}: \textit{Pantoea}, \textit{Pseudomonas}, \textit{Enterobacter}. Within a day all but the \textit{Enterobacteriaceae} are suppressed and lactic acid bacteria (LAB) take over
    \item The first LAB are not sourdough organisms: \textit{Leuconostoc}, \textit{Lactococcus}, \textit{Enterococcus}. At 30°C they fall below detection within a day or two
    \item Rye is the exception: it arrives with \textit{Weissella} already dominant
    \item Wink had a day-two starter analysed by two labs. The gas came from \textit{Leuconostoc citreum}, with a little \textit{Enterobacter cloacae}; neither lab found lactobacilli or yeast \caveat{not peer-reviewed}
    \item She links \textit{Enterobacter}, \textit{Aerococcus} and \textit{Leuconostoc} to the bad smell
  \end{itemize}
  \takeaway{The 3--4× rise and the smell on day 4 are this first wave. Keep to the schedule.}
  \source{Ercolini et al.\ (2013); Van der Meulen et al.\ (2007); Wink (2008b), not peer-reviewed}
\end{frame}

\begin{frame}{Day 5: the lull is acid shutting down the first wave}
  \begin{columns}[T]
    \begin{column}{0.66\textwidth}
      \small
      \begin{itemize}
        \item As acid builds, the first wave stops before the sourdough organisms have numbers: ``No significant growth!''
        \item Wink's pH paper: gassy bacteria grew at pH ≥\,5, not ≤\,4.5; yeast usually came the day after pH 3.5 \caveat{careful amateur work; no replicate counts}
        \item \hypothesis\ (hers): \textit{L.\ citreum}'s bacteriocins hold back the lactobacilli, so a bigger early rise means a longer lull
        \item Lab data agree: at 30°C, pH below 4.0 after day 2; yeast at 10⁷\,CFU/g after 4--5 days (one dough in four: 8)
      \end{itemize}
    \end{column}
    \begin{column}{0.31\textwidth}
      \centering
      \begin{tikzpicture}[x=1cm,y=1.05cm]
        \fill[crumb] (0,5) rectangle (0.5,6);
        \fill[crust!25] (0,3) rectangle (0.5,4.5);
        \draw[crust,thick] (0,3) rectangle (0.5,6);
        \foreach \p in {6,5,4.5,4,3.5,3}{
          \draw[crust] (0,\p) -- (-0.1,\p) node[left,font=\tiny,text=rye] {\p};
        }
        \node[font=\scriptsize\bfseries,text=crust] at (0.25,6.3) {pH};
        \node[right,font=\scriptsize,align=left,text width=2.6cm] at (0.6,5.5) {≥\,5: the gassy first wave grows};
        \node[right,font=\scriptsize,align=left,text width=2.6cm,text=muted] at (0.6,4.75) {between: the paper reads to 0.5};
        \node[right,font=\scriptsize,align=left,text width=2.6cm] at (0.6,4.15) {≤\,4.5: it does not};
        \draw[flag,very thick] (-0.05,3.5) -- (0.55,3.5);
        \node[right,font=\scriptsize,align=left,text width=2.6cm] at (0.6,3.5) {3.5: yeast usually the next day};
      \end{tikzpicture}
    \end{column}
  \end{columns}
  \takeaway{Almost no rise and a vomit smell on day 5 is normal. A pH strip at each feed, days 3--7, tests Wink's thresholds.}
  \source{Wink (2008a, b), not peer-reviewed; Van der Meulen et al.\ (2007)}
\end{frame}

\begin{frame}{Days 6--9: the sourdough bacteria take over, then little changes}
  \small
  Sourdough-specific LAB, such as \textit{Lactobacillus}, \textit{Pediococcus} and \textit{Weissella}, replace the first wave.
  \medskip

  \begin{tabularx}{\linewidth}{@{}l l >{\raggedright\arraybackslash}X@{}}
    \toprule
    \textbf{Where} & \textbf{Conditions} & \textbf{When it settles} \\
    \midrule
    Van der Meulen et al.\ (2007) & 30°C, refreshed daily & Well-adapted strains dominate from about day 5 \\
    Ercolini et al.\ (2013) & 25°C, 5\,h a day & Mature in 5--7 days; lactobacilli prominent by day 10 \\
    The video & 24°C & ``Bit of sour'' on day 7: about a day behind 30°C \\
    \bottomrule
  \end{tabularx}
  \takeaway{The first sour smell, with the rise back to 2×, is the hand-over: sourdough bacteria replacing the first wave.}
  \source{Van der Meulen et al.\ (2007); Ercolini et al.\ (2013); Full Proof Baking (2019)}
\end{frame}

\begin{frame}{Days 9--14: raising the ratio is about timing, not species}
  \small
  \begin{itemize}
    \item By day 9 the community has largely settled. What remains is matching food to appetite
    \item A 1:2:2 feed that peaks and collapses within 12\,h ran out of food well before the next feed, and then sits in its own acid
    \item A higher ratio dilutes the starter, so it takes longer to peak, and dilutes the acid it carried over: the same lever as chapter 11's recovery at 1:10:10
    \item ``Feed at peak'' is \textbf{a practitioner's rule, not a measured result}
    \item The plausible mechanism is chapter 11's compounding stressors: past peak the food is spent and acid and ethanol are at their highest, so a starter always fed after collapse is always fed at its most stressed
  \end{itemize}
  \takeaway{If it peaks well before the next feed is due, raise the ratio. That is the whole of days 9--13.}
  \source{Full Proof Baking (2019); chapter 11, \emph{Compounding Stressors} and \emph{Recovery Protocol}}
\end{frame}

% ===========================================================================
\section{Video, literature and this kitchen}

\begin{frame}{The method works; three of its explanations do not}
  \footnotesize
  \begin{tabularx}{\linewidth}{@{}c >{\raggedright\arraybackslash}p{3.9cm} X@{}}
    \toprule
    & \textbf{The video says} & \textbf{The evidence says} \\
    \midrule
    1 & It ``picks up the wild bacteria and yeasts from the air'' &
      \textbf{The flour, then the baker's hands.} Same recipe and flour on two continents: much of each community traced to the flour, the differences to the bakers' hands. In bakeries over a year, flour drove the LAB; the bakery's own microbes played a minor role. It happens in a covered jar too \\
    \addlinespace
    2 & Day 3: ``We've caught some wild yeast!'' &
      In a two-day-old starter the rise is bacterial. Her narration says ``not the correct balance'' at once: only the label is wrong \\
    \addlinespace
    3 & Day 8: a loose lid, as ``this is a live culture that requires air'' &
      Yeast and LAB ferment without oxygen. The lid is loose so CO₂ can escape, not so air can get in \\
    \bottomrule
  \end{tabularx}
  \takeaway{The cheesecloth does no harm, but it is not what inoculates the jar: the flour is.}
  \source{Reese et al.\ (2020); Minervini et al.\ (2018); Wink (2008a), not peer-reviewed}
\end{frame}

\begin{frame}{Start warm, not hot: 24°C or 30°C picks a different culture}
  \begin{columns}[T]
    \begin{column}{0.6\textwidth}
      \small
      \begin{itemize}
        \item New starters are small and unsettled, like chapter 11's depleted culture: 22--25°C, no selection pressure
        \item Warm is fine: her target is 27°C; Van der Meulen et al.\ (2007) had a stable culture in a week at 30°C
        \item \textit{L.\ citreum} is behind Wink's false start: at 24°C the first wave may never fully leave
        \item \hypothesis\ begun near 30°C, a starter may suit a warm-kept regime better (open tests 1--2)
      \end{itemize}
    \end{column}
    \begin{column}{0.37\textwidth}
      \centering
      {\scriptsize\color{muted} Daily backslopping: dominant LAB}\par\smallskip
      \begin{tikzpicture}
        \node[box=1.3cm] (t1) {\textbf{23°C}};
        \node[box=2.2cm,right=0.45cm of t1] (s1) {\textit{Leuconostoc citreum}};
        \node[box=1.3cm,below=0.35cm of t1] (t2) {\textbf{30°C, 37°C}};
        \node[box=2.2cm,right=0.45cm of t2] (s2) {\textit{Lactobacillus fermentum}};
        \draw[arrow] (t1) -- (s1);
        \draw[arrow] (t2) -- (s2);
      \end{tikzpicture}
      \par\smallskip
      {\scriptsize Van der Meulen et al.\ (2011), via chapter 09}
    \end{column}
  \end{columns}
  \takeaway{No hot builds until it passes the tripling test; then keep the mother at room temperature and build hot from it.}
  \source{Van der Meulen et al.\ (2007, 2011); Wink (2008b), not peer-reviewed; chapters 08, 09, 11}
\end{frame}

\begin{frame}{Tap water and ordinary flour will do; juice is optional}
  \begin{columns}[T]
    \begin{column}{0.47\textwidth}
      \small
      \textbf{\color{crust}Water}
      \begin{itemize}
        \item Bottled water is usually about chlorine. Danish tap water is groundwater, normally unchlorinated: tap water should do
      \end{itemize}
      \textbf{\color{crust}Flour}
      \begin{itemize}
        \item Any bread flour for the base; the Manitoba is more than strong enough
        \item Dark wholemeal rye (\textit{fuldkornsrugmel}): the whole grain is the point, as when chapter 11 adds wholemeal in recovery
      \end{itemize}
    \end{column}
    \begin{column}{0.5\textwidth}
      \small
      \textbf{\color{crust}A shortcut through the lull}
      \begin{itemize}
        \item Wink: pineapple juice (or orange juice, or apple cider) for water on days 1--3 acidifies the mix below what the gas producers tolerate
        \item Yeast on or before day 4, against about day 7 with water
        \item It ``doesn't change the end result''; never necessary, ``only patience''
        \item Here: juice for the day-1 mix and the day-3 feed \untested
      \end{itemize}
    \end{column}
  \end{columns}
  \source{Wink (2008a, b), not peer-reviewed. Juice in this method is the note's suggestion and has not been tried.}
\end{frame}

% ===========================================================================
\section{Drying a backup and reviving it}

\begin{frame}{A pinch worked where a soaked spoonful failed \ownresult}
  \begin{columns}[T]
    \begin{column}{0.48\textwidth}
      \begin{block}{First attempt}
        \small
        A recipe worked out in a September 2026 conversation with Claude:
        \begin{itemize}
          \item a larger amount of flakes
          \item soaked in warm water first
        \end{itemize}
        \textbf{It did not take.}
      \end{block}
    \end{column}
    \begin{column}{0.48\textwidth}
      \begin{block}{October 2026}
        \small
        Flakes a year old:
        \begin{itemize}
          \item a good pinch, about 1\,g, crushed into a fresh 30\,g water + 30\,g wholemeal/rye mix
          \item no separate soak
        \end{itemize}
        \textbf{Clearly active in several days.}
      \end{block}
    \end{column}
  \end{columns}
  \medskip
  {\small Why keep a backup: after a heat kill, it is the only way back to the original culture, rather than to whatever the flour supplies. That holds for the bacteria; the yeast may not come through drying.}
  \source{Own runs, September and October 2026; chapter 11, \emph{Recovery Protocol}}
\end{frame}

\begin{frame}{Dry it cool: a wet film is heat-fragile, a dry flake is not}
  \begin{columns}[T]
    \begin{column}{0.43\textwidth}
      \small
      \textbf{\color{crust}Drying it}
      \begin{itemize}
        \item Spread thin on baking paper; room temperature until it snaps, about 24--48\,h
        \item Crumble; store airtight, dry, cool and dark
        \item No oven, no warm dehydrator: the guide allows one only below 37°C
      \end{itemize}
    \end{column}
    \begin{column}{0.54\textwidth}
      \small
      \textbf{\color{crust}Dry is heat-proof; drying is not}
      \begin{itemize}
        \item Until dry, the film is a wet culture: six hours or more at 40--45°C is high risk, and the D-value is under a minute at 55°C (chapter 11)
        \item Dried \textit{S.\ cerevisiae} kept full viability through 60\,min at 75°C, and 10\,min at 100°C at 12\% relative humidity
        \item The September chat said dried cells lose heat tolerance: backwards
      \end{itemize}
    \end{column}
  \end{columns}
  \takeaway{The heat risk is the wet start of drying, not the dry end: dry at room temperature.}
  \source{Country Roads Sourdough, not peer-reviewed; Ribert et al.\ (2021); chapter 11}
\end{frame}

\begin{frame}{Reviving it in five steps, with amounts by eye \ownresult}
  \small
  \begin{enumerate}\setlength{\itemsep}{1pt}
    \item \textbf{Start:} crush a good pinch of flakes, about 1\,g, into about 30\,g water + 30\,g flour, half wholemeal, half rye. Water at room temperature or a little above, not cold; no separate soak
    \item \textbf{Wait, unfed, for several days} at 24--26°C until clearly active. A dry crust forms on top: the day-2 skin, harmless
    \item \textbf{Lift off the crust and feed what is underneath} with bread flour and wholemeal at about 1:1:1. From the next feed, keep 20\,g and discard the rest, as on day 4
    \item \textbf{Feed at peak} as in the main method; one or two feeds on the usual flour before baking with it
    \item \textbf{Dry a new backup} once it doubles reliably
  \end{enumerate}
  \takeaway{Record flake age, rough amount, water temperature, first sour smell, first rise, first doubling. Sour without a rise: the LAB are back, the yeast not yet.}
  \source{The method as run here, October 2026; amounts estimated, not weighed, and it does not seem sensitive to them}
\end{frame}

\begin{frame}{A dried starter is not a paused starter}
  \footnotesize
  \renewcommand{\arraystretch}{0.92}
  \begin{tabularx}{\linewidth}{@{}>{\raggedright\arraybackslash}p{3.0cm} >{\raggedright\arraybackslash}p{5.5cm} >{\raggedright\arraybackslash}X@{}}
    \toprule
    \textbf{What} & \textbf{Through drying} & \textbf{Why} \\
    \midrule
    Water & Leaves, to equilibrium with the room & --- \\
    CO₂ and ethanol & Leave & A gas and a volatile, from a thin film \\
    Acetic acid & About half leaves & Volatile: 57\% kept in freeze-drying \\
    Lactic acid & Mostly stays & Not volatile: 85\% kept in freeze-drying \\
    Spent flour: starch, gluten, bran, minerals & Stays & Already part-fermented: little food left in it \\
    Flour enzymes & Stay, idle until wet & Flour is itself a dry product whose enzymes work once wetted \\
    Cells & A minority survive dormant; most die & Next slide \\
    \bottomrule
  \end{tabularx}
  \takeaway{Slow air-drying probably loses at least as much acetic acid, so a flake brings back mostly lactic acid. That is why the dose matters.}
  \source{Acid figures: Kline (1979), a patent, not peer-reviewed}
\end{frame}

\begin{frame}{Drying kills most cells, the yeast hardest}
  \begin{columns}[T]
    \begin{column}{0.59\textwidth}
      \small
      \begin{itemize}
        \item First hours: the film keeps fermenting as sugars, salts and acids concentrate \caveat{reasoning; nothing measured}
        \item Damage: mechanical, structural and oxidative in yeast; in LAB mainly the membrane, as it changes phase
        \item Air-drying \textit{S.\ cerevisiae}: under 1 in a million growing cells survive; 20--40\% of cells that ran out of food
        \item \hypothesis\ so dry later than the usual ``at peak''? Uncertain: tolerance comes with respiration, a jar is short of oxygen, no sourdough yeast was tested (open test 4)
      \end{itemize}
    \end{column}
    \begin{column}{0.38\textwidth}
      \centering
      {\scriptsize\textbf{Survival, freeze-dried San Francisco starter}}\par\smallskip
      \begin{tikzpicture}[x=1.15cm,y=1cm]
        % log10 scale: x = log10(%) + 2, from 0.01 % to 100 %
        \foreach \x/\l in {0/0.01,1/0.1,2/1,3/10,4/100}{
          \draw[muted!40] (\x,-2.25) -- (\x,-2.35);
          \node[font=\tiny,text=muted] at (\x,-2.55) {\l\%};
        }
        \draw[muted!60] (0,-2.25) -- (4,-2.25);
        \foreach \y in {-0.45,-1.25}{\foreach \x in {1,2,3}{\draw[muted!30] (\x,\y-0.2) -- (\x,\y+0.2);}}
        \fill[own] (0,-0.25) rectangle (3.477,-0.65);
        \fill[own!45] (3.477,-0.25) rectangle (3.653,-0.65);
        \node[font=\tiny,anchor=south west,inner sep=1pt] at (0,-0.22) {\textit{L.\ sanfrancisco}, with additive: 30--45\%};
        \fill[own!60] (0,-1.05) rectangle (1.845,-1.45);
        \node[font=\tiny,anchor=south west,inner sep=1pt] at (0,-1.02) {\textit{L.\ sanfrancisco}, without: 0.7\%};
        \draw[flag,very thick] (0,-1.85) -- (0,-2.2);
        \node[font=\tiny,anchor=south west,inner sep=1pt] at (0,-1.82) {Yeast (\textit{Torulopsis holmii}): below 0.01\%};
      \end{tikzpicture}
      \par
      {\tiny Kline (1979) \caveat{a patent, not peer-reviewed}}\par
      {\tiny\color{muted} Freeze-drying is not air-drying, but this is the only sourdough-specific comparison found.}
    \end{column}
  \end{columns}
  \source{Calahan et al.\ (2011); Dupont et al.\ (2014); Santivarangkna et al.\ (2008); Kline (1979), a patent, not peer-reviewed}
\end{frame}


\begin{frame}{The bacteria come back on numbers, not toughness}
  \begin{columns}[T]
    \begin{column}{0.6\textwidth}
      \footnotesize
      \begin{itemize}\setlength{\itemsep}{1pt}
        \item Not toughness: without the additive, over 99\% of Kline's bacterium died
        \item Numbers: far more LAB than yeast (≈\,10⁷/g) to start, and a smaller share lost
        \item Why the bacterium dies less is not established (Potts 1994)
        \item Not the acid: acid-adapted \textit{L.\ lactis} dried no better; stopping growth is shown to protect only yeast
        \item Flour may protect a little \untested
      \end{itemize}
    \end{column}
    \begin{column}{0.37\textwidth}
      \centering
      {\scriptsize\textbf{Survivors in a 1\,g pinch, before storage}}\par\smallskip
      \begin{tikzpicture}[x=0.62cm,y=1cm]
        % log10 scale: x = log10(cells), from 1 to 10^6
        \foreach \x/\l in {0/1,2/10²,4/10⁴,6/10⁶}{
          \draw[muted!40] (\x,-1.85) -- (\x,-1.95);
          \node[font=\tiny,text=muted] at (\x,-2.15) {\l};
        }
        \draw[muted!60] (0,-1.85) -- (6,-1.85);
        \fill[own] (0,-0.25) rectangle (5.15,-0.65);
        \node[font=\tiny,anchor=south west,inner sep=1pt] at (0,-0.22) {LAB: more than 100\,000};
        \fill[flag!70] (0,-1.05) rectangle (3.3,-1.45);
        \node[font=\tiny,anchor=south west,inner sep=1pt] at (0,-1.02) {Yeast: fewer than ≈\,2\,000};
      \end{tikzpicture}
      \par\smallskip
      {\tiny\caveat{borrowed figures, not a count}}
    \end{column}
  \end{columns}
  \takeaway{Sour comes back first because the bacteria were never scarce; the rise waits on the yeast.}
  \source{Kline (1979), a patent, not peer-reviewed; Van der Meulen et al.\ (2007); Potts (1994); Calahan et al.\ (2011); Cnossen et al.\ (2020)}
\end{frame}

\begin{frame}{Water starts four things at once}
  \centering
  \begin{tikzpicture}[every node/.append style={font=\small,minimum height=1.45cm}]
    \node[box=6.5cm,align=left] (a) {\textbf{1. The lactic acid redissolves.} A pinch gives the fresh flour a small dose; a spoonful gives close to full-strength starter acid};
    \node[box=6.5cm,align=left,right=0.4cm of a] (b) {\textbf{2. Flour enzymes start freeing sugar again}, in the fresh flour and any left in the flake};
    \node[box=6.5cm,align=left,below=0.35cm of a] (c) {\textbf{3. Survivors take up water} and their membranes pass back through the phase change: the moment of leakage, worse in cold water};
    \node[box=6.5cm,align=left,below=0.35cm of b] (d) {\textbf{4. Dead cells leak} their contents into the mix};
    \path (a.north) -- (b.north) coordinate[midway] (m);
    \node[above=0.25cm of m,font=\small\bfseries,text=crust,minimum height=0pt] {Rehydration};
  \end{tikzpicture}
  \par\bigskip
  \raggedright\small
  Then a lag, while the survivors repair before they divide. From there the revival is a race between their regrowth and the flour's own microbes.
  \takeaway{Step 3 is where water temperature matters: the leakage is worse in cold water.}
  \source{Santivarangkna et al.\ (2008); Crowe et al.\ (1989)}
\end{frame}

\begin{frame}{A pinch skips the first wave; the yeast still has to build up}
  \centering
  \begin{tikzpicture}[node distance=0.13cm,every node/.append style={inner sep=2.4pt}]
    \node[font=\scriptsize\bfseries,text=crust,anchor=west] (ph) at (0,0) {From a pinch of dried flakes \ownresult};
    \node[cell=8.3cm,below=0.1cm of ph.south west,anchor=north west] (p0) {\textbf{Start:} about 1\,g crushed flakes + 30\,g water + 30\,g wholemeal/rye};
    \node[cell=8.3cm,below=of p0] (p1) {\textbf{First hours:} lactic acid redissolves · flour enzymes free sugar · survivors rehydrate and repair before dividing};
    \node[cell=8.3cm,below=of p1] (p2) {\textbf{Surviving LAB regrow:} many survivors, already acid-tolerant · pH falls below 4.5 → the flour's gassy first wave is shut out: no false rise, no lull};
    \node[cell=8.3cm,below=of p2] (p4) {\textbf{Yeast:} perhaps a hundred survivors, maybe none, + newcomers from flour, hands, kitchen · the acid admits the tolerant ones};
    \node[cell=8.3cm,below=of p4] (p5) {\textbf{Several days, unfed:} sour but no rise · dry crust on top · yeast regrows over about 20 doublings};
    \node[cell=8.3cm,below=of p5] (p6) {\textbf{First rise:} lift off the crust · feed bread flour + wholemeal 1:1:1 · then keep 20\,g, as on day 4};
    \node[cell=8.3cm,below=of p6] (p7) {Feed at peak until it doubles reliably → dry a new backup};
    \foreach \a/\b in {p0/p1,p1/p2,p2/p4,p4/p5,p5/p6,p6/p7}{\draw[arrow,thin] (\a) -- (\b);}
    % from scratch
    \node[font=\scriptsize\bfseries,text=crust,anchor=west] (sh) at (9.6,0) {From scratch, for comparison};
    \node[skipped=4.4cm,below=0.1cm of sh.south west,anchor=north west] (s1) {\textbf{Days 1--4:} the flour's first wave, \textit{Leuconostoc}, \textit{Enterobacter} · big rise, bad smell};
    \node[skipped=4.4cm,below=0.3cm of s1] (s2) {\textbf{Day 5: the lull} · acid stops the first wave};
    \node[cell=4.4cm,anchor=north west] (s3) at (s1.west |- p4.north) {\textbf{Days 6--9:} sourdough LAB take over · yeast the day after pH 3.5};
    \node[cell=4.4cm,anchor=north west] (s4) at (s1.west |- p6.north) {\textbf{Days 9--14:} ratios climb until it triples};
    \draw[arrow,thin] (s1) -- (s2);
    \draw[arrow,thin] (s2) -- (s3);
    \draw[arrow,thin] (s3) -- (s4);
    \draw[arrow,thin,dashed,draw=own] (p2.east) -- ++(0.35,0) |- ($(s1.south west)!0.5!(s2.north west)$)
      node[pos=0.75,above,font=\tiny,text=own] {skipped};
  \end{tikzpicture}
  \raggedright
  \source{own run, October 2026, for the timings. The order of events inside them is reasoning \caveat{no pH or counts were taken}}
\end{frame}

\begin{frame}{The surviving bacteria close the window the first wave needs}
  \small
  \begin{itemize}
    \item From scratch, a gassy first wave rises and collapses and the jar goes quiet before the sourdough organisms take over
    \item The revival took days, but had \textbf{no false start and no lull} \ownresult
    \item The likely reason is the inoculum: the surviving LAB are sourdough organisms, selected by months of feeding, that grow well at their own acidity
    \item They acidify the fresh flour before its first wave builds, closing Wink's window: the gas producers grow at pH ≥\,5 and not at ≤\,4.5
    \item From scratch nothing acid-tolerant is there in numbers, so the first wave gets days \caveat{reasoning that fits the run here; no pH was measured}
  \end{itemize}
  \takeaway{From flakes, expect quiet days that turn sour before rising: no day-4 surge, no lull.}
  \source{Wink (2008a), not peer-reviewed; own run, October 2026}
\end{frame}

\begin{frame}{The yeast restarts from perhaps a hundred cells, possibly none}
  \centering
  \begin{tikzpicture}[node distance=0.26cm,every node/.append style={nohyph}]
    \node[box=2.45cm,minimum height=1.35cm] (a) {\textbf{1\,g pinch}\\ about 2\,g of starter before drying};
    \node[box=2.45cm,minimum height=1.35cm,right=of a] (b) {\textbf{≈\,10⁷ yeast/g}\\ in a ripe starter, and far more LAB};
    \node[box=2.45cm,minimum height=1.35cm,right=of b] (c) {\textbf{Drying: under 0.01\%}\\ yeast survive; the bacterium 0.7\%, at least 70× better};
    \node[box=2.45cm,minimum height=1.35cm,right=of c] (d) {\textbf{A year stored: −94\%}\\ even for commercial dried yeast};
    \node[box=2.45cm,minimum height=1.35cm,right=of d,fill=own!12,draw=own] (e) {\textbf{Perhaps 100 cells}\\ possibly none: 1 or 2 per gram of fresh mix};
    \foreach \x/\y in {a/b,b/c,c/d,d/e}{\draw[arrow] (\x) -- (\y);}
    \node[font=\tiny,text=muted,below=0.05cm of b,minimum height=0pt] {Van der Meulen 2007};
    \node[font=\tiny,text=muted,below=0.05cm of c,minimum height=0pt] {Kline 1979};
    \node[font=\tiny,text=muted,below=0.05cm of d,minimum height=0pt] {Ribert 2021};
  \end{tikzpicture}
  \par\medskip\raggedright\small
  \begin{itemize}
    \item That is over 20 doublings short of a ripe starter: each tenfold is about 3.3 doublings
    \item The LAB start from far more survivors. The revival behaved as if the LAB had come through and the yeast had to be caught again \ownresult
    \item \caveat{arithmetic on borrowed figures, not a count}: Kline's starter was freeze-dried, and Ribert's yeast was grown to be dried
  \end{itemize}
  \source{Van der Meulen et al.\ (2007); Kline (1979), a patent, not peer-reviewed; Ribert et al.\ (2021)}
\end{frame}

\begin{frame}{It skipped the bacterial half of the first week, not the yeast half}
  \small
  \begin{itemize}
    \item \textbf{The yeast may not be the original.} Survivors compete with yeast arriving as it does from scratch: with the flour, and from the baker's hands
    \item A kitchen that kept the starter for years may hold its yeast: in four bakeries, the sourdough's dominant yeast also dominated walls, air, storage boxes and mixers. A ``new'' yeast may be the old one coming back \caveat{loose analogue}
    \item The yeast does not arrive in a from-scratch jar: the LAB have already acidified it. From scratch, yeast appeared the day after pH 3.5, so its clock starts at once, and the acid favours tolerant yeasts
  \end{itemize}
  \takeaway{A backup reliably brings back the bacteria, and perhaps the yeast. Drying later may save more of it (open test 4); whether it is the old yeast is open test 7.}
  \source{Reese et al.\ (2020); Minervini et al.\ (2015): bakeries hold far more flour and dough than a home, and the study cannot say which way the yeast moved; Wink (2008a), not peer-reviewed}
\end{frame}

\begin{frame}{A pinch, not a spoonful: the flakes carry the old starter's acid}
  \begin{columns}[T]
    \begin{column}{0.47\textwidth}
      \centering
      \begin{tikzpicture}[node distance=0.3cm,every node/.append style={nohyph}]
        \node[box=2.95cm] (sp) {\textbf{A spoonful}\\in a little water and flour};
        \node[box=2.95cm,below=of sp,minimum height=1.5cm] (spo) {Close to full-strength starter acid, little fresh food: chapter 11's dying culture};
        \node[box=2.95cm,right=0.3cm of sp] (pi) {\textbf{A pinch}\\in 60\,g fresh wholemeal and rye};
        \node[box=2.95cm,below=of pi,fill=own!12,draw=own,minimum height=1.5cm] (pio) {Far more dilute, plenty of food; wholegrain is what chapter 11 adds};
        \draw[arrow] (sp) -- (spo);
        \draw[arrow] (pi) -- (pio);
        \node[font=\scriptsize,below=0.1cm of spo,text=flag,align=center,text width=3.2cm] {first attempt, also warm-soaked: did not take \ownresult};
        \node[font=\scriptsize,below=0.1cm of pio,text=own] {worked \ownresult};
      \end{tikzpicture}
    \end{column}
    \begin{column}{0.5\textwidth}
      \small
      \begin{itemize}
        \item Flakes are the old flour plus all its lactic acid, which does not evaporate in drying
        \item Chapter 11 dilutes a dying culture out at 1:10:10
        \item Country Roads' recipe is the high-dose kind, 5\,g of flakes into 20\,g water + 20\,g flour, and expects about a week
        \item Two things changed, amount and soak: one run cannot separate them (open test 6). The soak is the less likely culprit: 35--38°C for half an hour is brief by chapter 11's standards, unless the water was hotter
      \end{itemize}
    \end{column}
  \end{columns}
  \source{Own runs, 2026; chapter 11, \emph{Recovery Protocol}; Country Roads Sourdough, not peer-reviewed}
\end{frame}

\begin{frame}{Water: not cold. The documented hazard is cold, not warm}
  \small
  \begin{itemize}
    \item The September chat advised lukewarm, 25--30°C, because cells are most fragile as water rushes back in. The fragility is real; the direction is wrong
    \item As a dry cell takes up water, its membrane lipids pass from gel to fluid. In cold water the membrane leaks while it changes state: \textbf{imbibitional damage}, measured in pollen, with dry yeast named as a case
    \item Active dry yeast rehydrated at 0°C against 40°C: less respiration and fermentation, membranes that did not reseal even after hours at 40°C, few cells growing on plates
    \item Limits: pollen; pure commercial \textit{S.\ cerevisiae}; flakes in flour may wet more slowly; nothing about 20°C. The direction is supported; the room-temperature penalty is unknown (test 5)
  \end{itemize}
  \takeaway{Mix with water at room temperature or a little above, up to about 35°C, then hold the jar at 24--26°C. Avoid cold tap water.}
  \source{Crowe et al.\ (1989); Peña et al.\ (1992); Country Roads Sourdough (room-temperature water; cold only slows things), not peer-reviewed}
\end{frame}

\begin{frame}{A backup loses ground every month, but slowly}
  \scriptsize
  \begin{tabularx}{\linewidth}{@{}l >{\raggedright\arraybackslash}X >{\raggedright\arraybackslash}p{4.6cm}@{}}
    \toprule
    \textbf{Source} & \textbf{What was stored} & \textbf{Loss} \\
    \midrule
    Calahan et al.\ (2011) & Air-dried \textit{S.\ cerevisiae} & About 20\% survival over a month \\
    Ribert et al.\ (2021) & Dried \textit{S.\ cerevisiae}, 25°C, 55\% relative humidity & 1.24 log, about 94\%, in a year \\
    Hajar-Azhari et al.\ (2024) & Spray-dried, kombucha-started sourdough powder, best of four packagings \caveat{loose analogue} & About 2.5 log of LAB and yeast in 90 days \\
    Here \ownresult & Flakes a year old (October 2026) & Back from a pinch in several days \\
    \bottomrule
  \end{tabularx}
  \smallskip
  \small
  \begin{itemize}
    \item A dry flake follows the room: water activity about 0.5 in a 50\% kitchen, far below anything that grows. Losses are mainly oxidation of membrane lipids: hence airtight, cool, dark
    \item Each tenfold loss costs about 3.3 doublings: an old backup wakes more slowly, not never. True of the LAB; the yeast is less certain
  \end{itemize}
  \takeaway{Replace the backup once a year from the live starter to keep the loss small.}
  \source{As in the table; FDA (1984); Santivarangkna et al.\ (2008); Country Roads Sourdough, not peer-reviewed; chapter 11}
\end{frame}

% ===========================================================================
\begin{frame}{Takeaways}
  \begin{itemize}
    \item The day-4 rise is bacteria and the day-5 lull is normal: keep to the schedule through both
    \item Feed at peak. When it falls before the next feed, go one ratio up
    \item Done means tripling in 4--6\,h on 1:2:2 at 27°C. No hot builds before that
    \item Start warm, not hot: 24°C or 30°C is a choice of culture as well as pace
    \item The microbes come from the flour and your hands, not the air
    \item Dry a backup thin at room temperature; store it airtight, cool and dark; replace it yearly
    \item Revive with a pinch in a fresh wholemeal and rye mix, water not cold. The bacteria come back first; the yeast later, and perhaps a new one
  \end{itemize}
\end{frame}

% ===========================================================================
\section{Open tests}

\begin{frame}{Seven open tests \untested}
  \scriptsize
  \begin{columns}[T]
    \begin{column}{0.45\textwidth}
      \textbf{\color{crust}Starting from scratch}
      \begin{enumerate}
        \item \textbf{Start temperature.} Two jars, same flour and day-1 mix: room temperature against 30°C (the chamber can hold it). Compare the end of the lull, the first tripling at 1:2:2, the smell. Van der Meulen et al.\ (2011) predicts different dominant LAB
        \item \textbf{Does a 30°C start take to the hot protocol better} than a 24°C one? Only once both pass the tripling test, on split lineages (chapter 08)
        \item \textbf{Juice versus water}, Wink's own design: two identical jars, one with juice on days 1--3
      \end{enumerate}
    \end{column}
    \begin{column}{0.53\textwidth}
      \textbf{\color{crust}The dried backup}
      \begin{enumerate}
        \setcounter{enumi}{3}
        \item \textbf{Dry at peak or after it.} Two films from one jar, at peak and a few hours after the fall; revive equal weights; compare the first doubling. Calahan et al.\ (2011) predicts the later film wins
        \item \textbf{Water temperature.} Same flakes, 35°C water against cold tap, both then at 24--26°C. Peña et al.\ (1992) predicts the cold jar is slower
        \item \textbf{Pinch versus spoonful}: a pinch against 10\,g, same flour and water temperature. The acid account predicts a slower spoonful; a pH strip at the start should show why
        \item \textbf{Same yeast or a new one?} No kitchen test; sequencing the fungal ITS region of flakes and revived starter, as Reese et al.\ (2020) did, would. Dead cells keep their DNA; a species match would not prove the same strain
      \end{enumerate}
    \end{column}
  \end{columns}
  \source{The Open tests section of notes/starter-from-scratch.md}
\end{frame}

% ===========================================================================
\begin{frame}{Sources}
  \tiny
  \begin{columns}[T]
    \begin{column}{0.49\textwidth}
      \textbf{\color{crust}The method and the first two weeks}\par\vspace{1.5pt}
      Kristen, Full Proof Baking (2019). \emph{How To Make A Sourdough Starter From Scratch.} YouTube, 7 Feb 2019.\par\vspace{1.5pt}
      Wink D (2008a). The pineapple juice solution, Part 1. \emph{Bread Lines} 16(1). Not peer-reviewed.\par\vspace{1.5pt}
      Wink D (2008b). The pineapple juice solution, Part 2. \emph{Bread Lines}, June 2008. Not peer-reviewed.\par\vspace{1.5pt}
      Van der Meulen R et al.\ (2007). \emph{Appl. Environ. Microbiol.} 73(15):4741--4750.\par\vspace{1.5pt}
      Van der Meulen R et al.\ (2011). \emph{Appl. Environ. Microbiol.} 77(8):2716--2726.\par\vspace{1.5pt}
      Ercolini D et al.\ (2013). Microbial ecology dynamics during rye and wheat sourdough preparation. \emph{Appl. Environ. Microbiol.} 79(24):7827--7836.\par\vspace{1.5pt}
      Minervini F et al.\ (2018). \emph{Frontiers in Microbiology} 9:1984.\par\vspace{1.5pt}
      Reese AT et al.\ (2020). Influences of ingredients and bakers on the bacteria and fungi in sourdough starters and bread. \emph{mSphere} 5(1):e00950-19.\par\vspace{1.5pt}
      Minervini F et al.\ (2015). House microbiotas as sources of LAB and yeasts in traditional Italian sourdoughs. \emph{Food Microbiology} 52:66--76.
    \end{column}
    \begin{column}{0.49\textwidth}
      \textbf{\color{crust}Drying and reviving}\par\vspace{1.5pt}
      Country Roads Sourdough. \emph{Dehydrating and Rehydrating Sourdough Starter: A How To Guide.} Not peer-reviewed.\par\vspace{1.5pt}
      Kline L (1979). Freeze-dried natural sour dough starter. US Patent 4,140,800. A patent, not peer-reviewed.\par\vspace{1.5pt}
      Calahan D et al.\ (2011). \emph{Genetics} 189(2):507--519.\par\vspace{1.5pt}
      Dupont S et al.\ (2014). \emph{Appl. Microbiol. Biotechnol.} 98(21):8821--8834.\par\vspace{1.5pt}
      Santivarangkna C et al.\ (2008). \emph{J. Appl. Microbiol.} 105(1):1--13.\par\vspace{1.5pt}
      Crowe JH et al.\ (1989). \emph{PNAS} 86(2):520--523.\par\vspace{1.5pt}
      Peña A et al.\ (1992). \emph{Arch. Microbiol.} 158(2):75--80.\par\vspace{1.5pt}
      Ribert P et al.\ (2021). \emph{Appl. Microbiol. Biotechnol.} 105(16--17):6409--6418.\par\vspace{1.5pt}
      Hajar-Azhari S et al.\ (2024). \emph{J. Food Sci. Technol.} 61(9):1701--1710.\par\vspace{1.5pt}
      Potts M (1994). \emph{Microbiol. Rev.} 58(4):755--805.\par\vspace{1.5pt}
      Cnossen DC et al.\ (2020). \emph{J. Dairy Res.} 87(3):382--385.\par\vspace{1.5pt}
      US FDA (1984). \emph{Water Activity (aw) in Foods.} Inspection Technical Guide.
    \end{column}
  \end{columns}
  \source{Full entries in the References section of notes/starter-from-scratch.md; chapter 11 of the main document}
\end{frame}

\end{document}
